\subsection{The hidden spread is a triangular spectrum}\label{tri:section}
The roots of $J_N$ have just collapsed to one point after division by
$e^\gamma\log N$.  That does not mean they have become equal.
It means we have chosen a scale on which their interesting differences
are hard to see.

A finer view reveals a familiar shape from an apparently distant
experiment.  Fill the region above a matrix diagonal with independent
complex Gaussian entries of variance $1/N$, and put zeros elsewhere:
\[
 \mathsf T_N=
 \begin{pmatrix}
 0& *& *&\cdots& *\\
 0& 0& *&\cdots& *\\
 \vdots&\vdots&\ddots&\ddots&\vdots\\
 0&0&\cdots&0&*\\
 0&0&\cdots&0&0
 \end{pmatrix}.
\]
All eigenvalues of this triangular matrix are zero.
Its singular values still measure how strongly it stretches vectors.
Their squares are the eigenvalues of $\mathsf T_N^*\mathsf T_N$;
these have a nontrivial limiting histogram.

That histogram is the Dykema--Haagerup law
$\operatorname{DH}_1$, supported on $[0,e]$
\cite{triDykemaHaagerup,triCheliotis}.
Its moments are
\[
 \int y^k\,d\operatorname{DH}_1(y)=\frac{k^k}{(k+1)!}
 \qquad(k=1,2,\ldots),
\]
so its mean is $1/2$.
No triangular matrices occur in the definition of $J_N$.
Nevertheless its roots reproduce this same law.

\begin{theorem}[The triangular spectrum]\label{tri:main}
Write $-R_{N,1},\ldots,-R_{N,N}$ for the roots of $J_N$,
with multiplicity. Center their positive magnitudes by setting
\[
 t_{N,j}=e^{-\gamma}R_{N,j}-\log N-\log\log N.
\]
Then
\begin{equation}\label{tri:exp}
 \boxed{\frac1N\sum_{j=1}^N
 \delta_{\exp(t_{N,j})}
 \Longrightarrow\operatorname{DH}_1.}
\end{equation}
Equivalently, if $T$ is the logarithm of a random variable with law
$\operatorname{DH}_1$, then
\begin{equation}\label{tri:centered}
 \frac1N\sum_{j=1}^N
 \delta_{t_{N,j}}
 \Longrightarrow\mathcal L(T).
\end{equation}
The law of $T$ has support $(-\infty,1]$; its explicit density is
given in \eqref{tri:density}.
\end{theorem}

Here is a practical reading of the theorem.
List the positive magnitudes of the arithmetic roots.
Multiply them by $e^{-\gamma}$, subtract
$\log N+\log\log N$, and make a histogram.
The result approaches a fixed asymmetric shape.
Exponentiate the centered entries, and the shape becomes the
squared-singular-value histogram of triangular Gaussian noise.

The common centering matters.
Dividing by $\log N$ erases fluctuations on a constant scale;
subtracting both $\log N$ and $\log\log N$ exposes them.
The extra $\log\log N$ is part of the arithmetic calculation.

\paragraph{Where the bridge comes from.}
The proof has three steps, all supplied in Appendix~\ref{triapp:all}.
First form the entire function
\[
 \mathcal E(z)=\sum_{k\ge0}K(k)\frac{z^k}{k!}.
\]
Its zeros are negative.  If their magnitudes are $\rho_j$, the prime
number theorem leads to the counting law
\[
 \#\{j:\rho_j\le r\}\sim e^{-\gamma}\frac r{\log^2r}.
\]
Second, Campbell and Jalowy's transfer theorem
\cite[Corollary 2.3]{finCampbellJalowy} turns this growth of the
entire function's zeros into a limiting law for the polynomial roots.
Finally, the precise arithmetic centering identifies that law as the
logarithm of $\operatorname{DH}_1$, using the identification of
Arizmendi and Hasebe \cite[Proposition 4.4]{triArizmendiHasebe}.

The connection therefore has a definite route:
coprimality probabilities determine an entire function;
its supply of zeros determines a free-stable law;
exponentiating that law gives the triangular spectrum.
The external transfer theorem supplies the bridge.
The prime estimates determine which law and which scale cross it.

The histogram alone does not locate the largest arithmetic root.
Question~\ref{tri:edgequestion} asks whether it also reaches the edge
of the triangular spectrum.
